%%
% The BIThesis Template for Undergraduate Thesis
%
% 北京理工大学毕业设计（论文） —— 使用 XeLaTeX 编译
%
% Copyright 2021-2023 BITNP
%
% This work may be distributed and/or modified under the
% conditions of the LaTeX Project Public License, either version 1.3
% of this license or (at your option) any later version.
% The latest version of this license is in
%   http://www.latex-project.org/lppl.txt
% and version 1.3 or later is part of all distributions of LaTeX
% version 2005/12/01 or later.
%
% This work has the LPPL maintenance status `maintained'.
%
% The Current Maintainer of this work is Feng Kaiyu.
%
% Compile with: xelatex -> biber -> xelatex -> xelatex
%%

% !TeX program = xelatex
% !BIB program = biber

\documentclass[type=bachelor_english]{bithesis}

\BITSetup{
  cover = {
    % 封面需要「北京理工大学」字样图片，如无必要请勿修改该项。
    headerImage = images/header.png,
    % 封面标题需要“华文细黑”，如无必要请勿修改该项。
    xiheiFont = STXIHEI.TTF,
    % 本科生盲审要求删去封面，而不是隐藏封面信息。
    hideCoverInPeerReview = true,
  },
  info = {
    % 想要删除某项封面信息，直接删除该项即可。
    % 想要让某项封面信息留空（但是保留下划线），请设置为空字符串。
    % 如需要换行，则用 “\\” 符号分割。
    titleEn = Your Thesis Title,
    school = School of Mechanical Engineering,
    major = Bechelor of Science in Mechanical Engineering,
    author = Feng Kaiyu,
    studentId = 11xxxxxxxx,
    supervisor = Alex Zhang,
    keywordsEn = {Computer-Aided Design; FEM; CAM},
    % 如果你的毕设为校外毕设，请将下面这一行语句解除注释（删除第一个百分号字符）并填写你的校外毕设导师名字
    % externalSupervisor = 左偏树,
  },
  style = {
    % 保持参考文献的缩进样式与 Word 模板一致。
    % 如果你不需要此样式，请将此行注释掉。
    bibliographyIndent = false,
    % 如无必要请勿修改该项。
    % head = {自定义页眉文字}
  }
}

% 使用 listings 宏包进行代码块使用，并使用了预定义的样式，
% 你也可以选用自己的喜欢的其他宏包，如 minted；
% 然而由于 minted 依赖 Python 的 Pygments 库作为外部依赖，因此出于模板的建议性考虑，我们没有提供 minted 进行代码块书写的示例。
% 但是，我们仍旧非常建议你使用 minted。
\usepackage{listings}

\usepackage[
  backend=biber,
  style=gb7714-2015,
  gbalign=gb7714-2015,
  gbnamefmt=lowercase,
  gbpub=false,
  doi=false,
  url=false,
  eprint=false,
  isbn=false,
]{biblatex}

% 参考文献引用文件位于 misc/ref.bib
\addbibresource{misc/ref.bib}


% 文档开始
\begin{document}

% 标题页面：如无特殊需要，本部分无需改动
% %%
% The BIThesis Template for Bachelor Graduation Thesis
%
% 北京理工大学毕业设计（论文）封面页 —— 使用 XeLaTeX 编译
%
% Copyright 2020-2022 BITNP
%
% This work may be distributed and/or modified under the
% conditions of the LaTeX Project Public License, either version 1.3
% of this license or (at your option) any later version.
% The latest version of this license is in
%   http://www.latex-project.org/lppl.txt
% and version 1.3 or later is part of all distributions of LaTeX
% version 2005/12/01 or later.
%
% This work has the LPPL maintenance status `maintained'.
%
% The Current Maintainer of this work is Feng Kaiyu.
%
% 封面
%
% 如无特殊需要，本页面无需更改

% Underline new command for student information
% Usage: \dunderline[<offset>]{<line_thickness>}
\newcommand\dunderline[3][-1pt]{{%
  \setbox0=\hbox{#3}
  \ooalign{\copy0\cr\rule[\dimexpr#1-#2\relax]{\wd0}{#2}}}}

% Cover Page
\begin{titlepage}
  \makeatletter
  \@ifundefined{externalMentorName}{
    % 校内毕设封面顶部间距
    \vspace*{19mm}
  }{
    % 校外毕设封面顶部间距
    \vspace*{13mm}
  }
  \centering

  \includegraphics[width=9.87cm]{images/header.png}

  \vspace*{-3mm}

  \zihao{-0}\textbf{\ziju{0.12}\songti{本科生毕业设计（论文）}}

  \vspace{16mm}

  \zihao{2}\textbf{\xihei\thesisTitle}

  \vspace{3mm}

  \begin{spacing}{1.2}
    \zihao{3}\selectfont{\textbf{\thesisTitleEN}}
  \end{spacing}

  \vspace{15mm}

  \flushleft

  \makeatletter
  \@ifundefined{externalMentorName}{
    % 生成校内毕设封面字段
    \makeatother
    \begin{spacing}{1.8}
      \hspace{27mm}\songti\zihao{3}\selectfont{学\hspace{11mm}院：\dunderline[-10pt]{1pt}{\makebox[78mm][c]{\deptName}}}

      \hspace{27mm}\songti\zihao{3}\selectfont{专\hspace{11mm}业：\dunderline[-10pt]{1pt}{\makebox[78mm][c]{\majorName}}}

      \hspace{27mm}\songti\zihao{3}\selectfont{学生姓名：\dunderline[-10pt]{1pt}{\makebox[78mm][c]{\yourName}}}

      \hspace{27mm}\songti\zihao{3}\selectfont{学\hspace{11mm}号：\dunderline[-10pt]{1pt}{\makebox[78mm][c]{\yourStudentID}}}

      \hspace{27mm}\songti\zihao{3}\selectfont{指导教师：\dunderline[-10pt]{1pt}{\makebox[78mm][c]{\mentorName}}}
    \end{spacing}
  }{
    % 生成校外毕设封面字段
    \makeatother
    \begin{spacing}{1.8}
      \hspace{19.4mm}\songti\zihao{3}\selectfont{学\hspace{19.6mm}院\hspace{3mm}：\dunderline[-10pt]{1pt}{\makebox[77.4mm][c]{\deptName}}}

      \hspace{19.4mm}\songti\zihao{3}\selectfont{专\hspace{19.6mm}业\hspace{3mm}：\dunderline[-10pt]{1pt}{\makebox[77.4mm][c]{\majorName}}}

      \hspace{19.4mm}\songti\zihao{3}\selectfont{学\hspace{2.8mm}生\hspace{2.8mm}姓\hspace{2.8mm}名\hspace{3mm}：\dunderline[-10pt]{1pt}{\makebox[77.4mm][c]{\yourName}}}

      \hspace{19.4mm}\songti\zihao{3}\selectfont{学\hspace{19.6mm}号\hspace{3mm}：\dunderline[-10pt]{1pt}{\makebox[77.4mm][c]{\yourStudentID}}}

      \hspace{19.4mm}\songti\zihao{3}\selectfont{指\hspace{2.8mm}导\hspace{2.8mm}教\hspace{2.8mm}师\hspace{3mm}：\dunderline[-10pt]{1pt}{\makebox[77.4mm][c]{\mentorName}}}

      \hspace{19.4mm}\songti\zihao{3}\selectfont{校外指导教师：\dunderline[-10pt]{1pt}{\makebox[77.4mm][c]{\externalMentorName}}}
    \end{spacing}
  }

  \vspace*{\fill}
  \centering
  \zihao{3}\ziju{0.5}\songti{\today}
\end{titlepage}

\MakeCover

% 前置页面定义
\frontmatter

\MakeOriginality

% 摘要：在摘要相应的 TeX 文件处进行摘要部分的撰写
%%
% The BIThesis Template for Bachelor Graduation Thesis
%
% 北京理工大学毕业设计（论文）中英文摘要 —— 使用 XeLaTeX 编译
%
% Copyright 2020 Spencer Woo
%
% This work may be distributed and/or modified under the
% conditions of the LaTeX Project Public License, either version 1.3
% of this license or (at your option) any later version.
% The latest version of this license is in
%   http://www.latex-project.org/lppl.txt
% and version 1.3 or later is part of all distributions of LaTeX
% version 2005/12/01 or later.
%
% This work has the LPPL maintenance status `maintained'.
%
% The Current Maintainer of this work is Spencer Woo.

% 中英文摘要章节
\topskip=0pt
\zihao{-4}

\vspace*{-6mm}

\begin{center}
  \heiti\zihao{-2}\textbf{\thesisTitle}
\end{center}

\vspace*{2mm}

\addcontentsline{toc}{chapter}{摘~~~~要}
{\let\clearpage\relax \chapter*{摘~~~~要}}
\setcounter{page}{1}

\vspace*{1mm}

\setstretch{1.53}
\setlength{\parskip}{0em}

% 中文摘要正文
本文……。

\textcolor{blue}{摘要正文选用模板中的样式所定义的“正文”，每段落首行缩进 2 个字符；或者手动设置成每段落首行缩进 2 个汉字，字体：宋体，字号：小四，行距：固定值 22 磅，间距：段前、段后均为 0 行。阅后删除此段。}

\textcolor{blue}{摘要是一篇具有独立性和完整性的短文，应概括而扼要地反映出本论文的主要内容。包括研究目的、研究方法、研究结果和结论等，特别要突出研究结果和结论。中文摘要力求语言精炼准确，本科生毕业设计（论文）摘要建议 300-500 字。摘要中不可出现参考文献、图、表、化学结构式、非公知公用的符号和术语。英文摘要与中文摘要的内容应一致。阅后删除此段。}

\vspace{4ex}\noindent\textbf{\heiti 关键词：北京理工大学；本科生；毕业设计（论文）}
\newpage

% 英文摘要章节
\topskip=0pt

\vspace*{2mm}

\begin{spacing}{0.95}
  \centering
  \heiti\zihao{3}\textbf{\thesisTitleEN}
\end{spacing}

\vspace*{16mm}

\addcontentsline{toc}{chapter}{Abstract}
{\let\clearpage\relax \chapter*{Abstract}}
\setcounter{page}{2}

\setstretch{1.53}
\setlength{\parskip}{0em}
In order to study……

\textcolor{blue}{Abstract 正文设置成每段落首行缩进 2 字符，字体：Times New Roman，字号：小四，行距：固定值 22 磅，间距：段前、段后均为 0 行。阅后删除此段。}

\vspace{3ex}\noindent\textbf{Key Words: BIT； Undergraduate； Graduation Project (Thesis)}
\newpage


\MakeTOC

% 正文开始
\mainmatter

% 第一章
%%
% The BIThesis Template for Bachelor Graduation Thesis
%
% 北京理工大学毕业设计（论文）第一章节 —— 使用 XeLaTeX 编译
%
% Copyright 2020-2022 BITNP
%
% This work may be distributed and/or modified under the
% conditions of the LaTeX Project Public License, either version 1.3
% of this license or (at your option) any later version.
% The latest version of this license is in
%   http://www.latex-project.org/lppl.txt
% and version 1.3 or later is part of all distributions of LaTeX
% version 2005/12/01 or later.
%
% This work has the LPPL maintenance status `maintained'.
%
% The Current Maintainer of this work is Feng Kaiyu.
%
% 第一章节

\chapter{一级题目}

\section{二级题目}
% 这里插入一个参考文献，仅作参考

\subsection{三级题目}

正文……\cite{yuFeiJiZongTiDuoXueKeSheJiYouHuaDeXianZhuangYuFaZhanFangXiang2008}……\cite{Hajela2012Application}

\textcolor{blue}{正文部分：宋体、小四；正文行距：22磅；间距段前段后均为0行。阅后删除此段。}

\textcolor{blue}{图、表居中，图注标在图下方，表头标在表上方，宋体、五号、居中，1.25倍行距，间距段前段后均为0行，图表与上下文之间各空一行。阅后删除此段。}

\textcolor{blue}{\underline{\underline{图-示例：（阅后删除此段）}}}


\begin{figure}[htbp]
  \centering
  \includegraphics[]{images/bit_logo.png}
  \caption{标题序号}\label{标题序号} % label 用来在文中索引
\end{figure}

\textcolor{blue}{\underline{\underline{表-示例：（阅后删除此段）}}}
% 三线表
\begin{table}[htbp]
  \linespread{1.5}
  \zihao{5}
  \centering
  \caption{统计表}\label{统计表}
  \begin{tabular}{*{5}{>{\centering\arraybackslash}p{2cm}}} \toprule
    项目    & 产量    & 销量    & 产值   & 比重    \\ \hline
    手机    & 1000  & 10000 & 500  & 50\%  \\
    计算机   & 5500  & 5000  & 220  & 22\%  \\
    笔记本电脑 & 1100  & 1000  & 280  & 28\%  \\
    合计    & 17600 & 16000 & 1000 & 100\% \\ \bottomrule
    \end{tabular}
\end{table}

\textcolor{blue}{公式标注应于该公式所在行的最右侧。对于较长的公式只可在符号处（+、-、*、/、$\leqslant$ $\geqslant$ 等）转行。在文中引用公式时，在标号前加“式”，如式（1-2）。阅后删除此
段。}

\textcolor{blue}{公式-示例：（阅后删除此段）}
% 公式上下不要空行，置于同一个段落下即可，否则上下距离会出现高度不一致的问题
\begin{equation}
    LRI=1\ ∕\ \sqrt{1+{\left(\frac{{\mu }_{R}}{{\mu }_{s}}\right)}^{2}{\left(\frac{{\delta }_{R}}{{\delta }_{s}}\right)}^{2}}
\end{equation}

\subsubsection{生僻字}

% 一个可能无法正常显示的生僻字
一个可能无法正常显示的生僻字: 彧。下文注释中，介绍了如何通过自定义字体来显示生僻字。

% 定义一个提供了生僻字的字体，注意要确保你的系统存在该字体
% \setCJKfamilyfont{custom-font}{Noto Serif CJK SC}

% 使用自己定义的字体
% 使用提供了相应字型的字体：\CJKfamily{custom-font}{彧}。


% 在这里添加第二章、第三章的 TeX 文件的引用
%%
% The BIThesis Template for Bachelor Graduation Thesis
%
% 北京理工大学毕业设计（论文）第二章节 —— 使用 XeLaTeX 编译
%
% Copyright 2020-2021 BITNP
%
% This work may be distributed and/or modified under the
% conditions of the LaTeX Project Public License, either version 1.3
% of this license or (at your option) any later version.
% The latest version of this license is in
%   http://www.latex-project.org/lppl.txt
% and version 1.3 or later is part of all distributions of LaTeX
% version 2005/12/01 or later.
%
% This work has the LPPL maintenance status `maintained'.
%
% The Current Maintainer of this work is Feng Kaiyu.
%%

\chapter{另一个章节}

\section{代码片段}

\begin{lstlisting}[language=Python, caption={Python Code}, label={lst:pythonfile}]
import numpy as np

def incmatrix(genl1,genl2):
    m = len(genl1)
    n = len(genl2)
    M = None #to become the incidence matrix
    VT = np.zeros((n*m,1), int)  #dummy variable

    #compute the bitwise xor matrix
    M1 = bitxormatrix(genl1)
    M2 = np.triu(bitxormatrix(genl2),1)

    for i in range(m-1):
        for j in range(i+1, m):
            [r,c] = np.where(M2 == M1[i,j])
            for k in range(len(r)):
                VT[(i)*n + r[k]] = 1;
                VT[(i)*n + c[k]] = 1;
                VT[(j)*n + r[k]] = 1;
                VT[(j)*n + c[k]] = 1;

                if M is None:
                    M = np.copy(VT)
                else:
                    M = np.concatenate((M, VT), 1)

                VT = np.zeros((n*m,1), int)

    return M
\end{lstlisting}

\chapter{Engineering Design}

Although engineering drawing still plays an important role in product design and manufacturing in many industrial sectors around the world, manual sketching for creating drawings has been gradually replaced by CAD (computer-aided design) software using computers. Beginning in the 1980s, CAD software reduced the need for draftsmen significantly, especially in small to mid-sized companies. The software’s affordability and ability to run on personal computers in the mid-1990s allowed engineers to do their own drafting and analytic work to some extent \ref{eq:1}.

\begin{equation}
x^n + y^n = z^n
\label{eq:1}
\end{equation}


% 后置内容
\backmatter

% 结论：在结论相应的 TeX 文件处进行结论部分的撰写
%%
% The BIThesis Template for Undergraduate Thesis
%
% 北京理工大学毕业设计（论文） —— 使用 XeLaTeX 编译
%
% Copyright 2021-2022 BITNP
%
% This work may be distributed and/or modified under the
% conditions of the LaTeX Project Public License, either version 1.3
% of this license or (at your option) any later version.
% The latest version of this license is in
%   http://www.latex-project.org/lppl.txt
% and version 1.3 or later is part of all distributions of LaTeX
% version 2005/12/01 or later.
%
% This work has the LPPL maintenance status `maintained'.
%
% The Current Maintainer of this work is Feng Kaiyu.
%
% Compile with: xelatex -> biber -> xelatex -> xelatex
%%

\begin{conclusion}
  During the 1980s, design-for-manufacturing practices were put into place in thousands of firms. Today DFM is an essential part of almost every product development effort. No longer can designers “throw the design over the wall” to production engineers. As a result of this emphasis on improved design quality, some manufacturers claim to have reduced production costs of products by up to 50 percent. In fact, comparing current new product designs with earlier generations, one can usually identify fewer parts in the new product, as well as new materials, more integrated and custom parts, higher-volume standard parts and subassemblies, and simpler assembly procedures.
\end{conclusion}

% 参考文献：如无特殊需要，参考文献相应的 TeX 文件无需改动，添加参考文献请使用 BibTeX 的格式
%   添加至 misc/ref.bib 中，并在正文的相应位置使用 \cite{xxj} 的格式引用参考文献
%%
% The BIThesis Template for Undergraduate Thesis
%
% 北京理工大学毕业设计（论文） —— 使用 XeLaTeX 编译
%
% Copyright 2021-2022 BITNP
%
% This work may be distributed and/or modified under the
% conditions of the LaTeX Project Public License, either version 1.3
% of this license or (at your option) any later version.
% The latest version of this license is in
%   http://www.latex-project.org/lppl.txt
% and version 1.3 or later is part of all distributions of LaTeX
% version 2005/12/01 or later.
%
% This work has the LPPL maintenance status `maintained'.
%
% The Current Maintainer of this work is Feng Kaiyu.
%
% Compile with: xelatex -> biber -> xelatex -> xelatex
%%

\begin{bibprint}
  \printbibliography[heading=none]
\end{bibprint}

% 附录：在附录相应的 TeX 文件处进行附录部分的撰写
%%
% The BIThesis Template for Graduate Thesis
%
% Copyright 2020-2023 Yang Yating, BITNP
%
% This work may be distributed and/or modified under the
% conditions of the LaTeX Project Public License, either version 1.3
% of this license or (at your option) any later version.
% The latest version of this license is in
%   http://www.latex-project.org/lppl.txt
% and version 1.3 or later is part of all distributions of LaTeX
% version 2005/12/01 or later.
%
% This work has the LPPL maintenance status `maintained'.
%
% The Current Maintainer of this work is Feng Kaiyu.

\begin{appendices}
  \chapter{费马大定理的证明}
  关于此，我确信已发现了一种美妙的证法，可惜这里空白的地方太小，写不下。

  \chapter{Maxwell Equations}
  因为在柱坐标系下，$\overline{\overline\mu}$是对角的，所以Maxwell方程组中电场$\bf
  E$的旋度

  所以$\bf H$的各个分量可以写为：
  \begin{subequations}
    \begin{eqnarray}
      H_r=\frac{1}{\mathbf{i}\omega\mu_r}\frac{1}{r}\frac{\partial
        E_z}{\partial\theta } \\
      H_\theta=-\frac{1}{\mathbf{i}\omega\mu_\theta}\frac{\partial E_z}{\partial r}
    \end{eqnarray}
  \end{subequations}

  同样地，在柱坐标系下，$\overline{\overline\epsilon}$是对角的，所以Maxwell方程组中磁场$\bf
  H$的旋度
  \begin{subequations}
    \begin{eqnarray}
      &&\nabla\times{\bf H}=-\mathbf{i}\omega{\bf D}\\
      &&\left[\frac{1}{r}\frac{\partial}{\partial
          r}(rH_\theta)-\frac{1}{r}\frac{\partial
          H_r}{\partial\theta}\right]{\hat{\bf
          z}}=-\mathbf{i}\omega{\overline{\overline\epsilon}}{\bf
        E}=-\mathbf{i}\omega\epsilon_zE_z{\hat{\bf z}} \\
      &&\frac{1}{r}\frac{\partial}{\partial
        r}(rH_\theta)-\frac{1}{r}\frac{\partial
        H_r}{\partial\theta}=-\mathbf{i}\omega\epsilon_zE_z
    \end{eqnarray}
  \end{subequations}

  由此我们可以得到关于$E_z$的波函数方程：
  \begin{eqnarray}
    \frac{1}{\mu_\theta\epsilon_z}\frac{1}{r}\frac{\partial}{\partial r}
    \left(r\frac{\partial E_z}{\partial r}\right)+
    \frac{1}{\mu_r\epsilon_z}\frac{1}{r^2}\frac{\partial^2E_z}{\partial\theta^2}
    +\omega^2 E_z=0
  \end{eqnarray}

  \chapter{要求}

  \textcolor{blue}{
  有些材料编入文章主体会有损于编排的条理性和逻辑性，或有碍于文章结构的紧凑和突出主题思想等，这些材料可作为附录另页排在参考文献之后，也可以单编成册。下列内容可作为附录：
  }
  \begin{enumerate}
    \item \textcolor{blue}{为了整篇论文材料的完整，但编入正文有损于编排的条理性和逻辑性的材料，这一类材料包括比正文更为详尽的信息、研究方法和技术等更深入的叙述，以及建议可阅读的参考文献题录和对了解正文内容有用的补充信息等；}
    \item \textcolor{blue}{ 由于篇幅过大或取材的复制资料不便于编入正文的材料； }
    \item \textcolor{blue}{ 不便于编入正文的罕见珍贵资料； }
    \item \textcolor{blue}{ 一般读者无须阅读，但对本专业同行有参考价值的资料； }
    \item \textcolor{blue}{ 某些重要的原始数据、推导、计算程序、框图、结构图、注释、统计表、计算机打印输出件等； }
  \end{enumerate}

  \section{一级标题}
  \subsection{二级标题}
\end{appendices}

% 致谢：在致谢相应的 TeX 文件处进行致谢部分的撰写
%%
% The BIThesis Template for Undergraduate Thesis
%
% 北京理工大学毕业设计（论文） —— 使用 XeLaTeX 编译
%
% Copyright 2021-2023 BITNP
%
% This work may be distributed and/or modified under the
% conditions of the LaTeX Project Public License, either version 1.3
% of this license or (at your option) any later version.
% The latest version of this license is in
%   http://www.latex-project.org/lppl.txt
% and version 1.3 or later is part of all distributions of LaTeX
% version 2005/12/01 or later.
%
% This work has the LPPL maintenance status `maintained'.
%
% The Current Maintainer of this work is Feng Kaiyu.
%
% Compile with: xelatex -> biber -> xelatex -> xelatex
%%

\begin{acknowledgements}
acknowledgements  This book contains material developed for use 
in the interdisciplinary courses on product development that we teach. 
Participants in these courses include graduate students in engineering, 
industrial design students, and MBA students. While we 
aimed the book at interdisciplinary graduate-level audiences such 
as this, many faculty teaching graduate and undergraduate courses 
in engineering design have also found the material useful. Product 
Design and Development is also for practicing professionals. Indeed, 
we could not avoid writing  for  a  professional  audience,  
because  most  of  our  students  are  themselvesprofessionals who have worked 
either in product development or in closely related functions.

This book blends the perspectives of marketing, design, and manufacturing into 
a single approach to product development. As a result, we provide students 
of all kinds with an appreciation for the realities of industrial practice 
and for the complex and essential roles played by the various members of
product development teams. For industrial practitioners, in  particular,  
we  provide  a  set  of  product  development  methods  that  can  be  
put  into immediate practice on development projects.
\end{acknowledgements}


\end{document}
